% 正文片段。主文档导言区需加入：\usepackage{pgfplots}
% 并设置：\pgfplotsset{compat=1.18}
% 建议使用 XeLaTeX 编译中文标签。
\subsection{分行业原始 PE}

图~\ref{fig:industry-raw-pe} 使用 2025 年 12 月 31 日最新截面，按申万一级行业汇总
未经 MAD 去极值和 Z 标准化的原始 PE。为避免接近零的 EP 使 PE 均值被极端值主导，
柱高采用行业内原始 PE 的中位数。亏损公司的负 PE 仍保留在统计中。

\begin{figure}[p]
  \centering
  \pgfplotsset{compat=1.18}
  \begin{tikzpicture}
    \begin{axis}[
      xbar,
      width=0.94\textwidth,
      height=0.80\textheight,
      xmin=-8,
      xmax=68,
      bar width=4.2pt,
      symbolic y coords={
        房地产,
    银行,
    建筑装饰,
    传媒,
    煤炭,
    纺织服饰,
    钢铁,
    交通运输,
    建筑材料,
    环保,
    公用事业,
    非银金融,
    石油石化,
    轻工制造,
    农林牧渔,
    医药生物,
    食品饮料,
    家用电器,
    基础化工,
    计算机,
    社会服务,
    电力设备,
    美容护理,
    汽车,
    商贸零售,
    机械设备,
    有色金属,
    电子,
    国防军工,
    综合,
    通信
      },
      ytick={
        房地产,
    银行,
    建筑装饰,
    传媒,
    煤炭,
    纺织服饰,
    钢铁,
    交通运输,
    建筑材料,
    环保,
    公用事业,
    非银金融,
    石油石化,
    轻工制造,
    农林牧渔,
    医药生物,
    食品饮料,
    家用电器,
    基础化工,
    计算机,
    社会服务,
    电力设备,
    美容护理,
    汽车,
    商贸零售,
    机械设备,
    有色金属,
    电子,
    国防军工,
    综合,
    通信
      },
      yticklabel style={font=\scriptsize},
      xlabel={行业内原始 PE 中位数（倍）},
      xlabel style={font=\small},
      xtick distance=10,
      xmajorgrids=true,
      grid style={gray!25},
      axis line style={gray!65},
      tick style={gray!65},
      enlarge y limits=0.012,
      point meta=x,
      nodes near coords={\pgfmathprintnumber[fixed,precision=2]{\pgfplotspointmeta}},
      every node near coord/.append style={font=\tiny},
    ]
      \addplot+[
        draw=none,
        fill={rgb,255:red,59;green,111;blue,182},
        bar shift=0pt
      ] coordinates {
        (6.0471,银行)
    (11.6142,建筑装饰)
    (12.2231,传媒)
    (13.4089,煤炭)
    (15.9799,纺织服饰)
    (16.3224,钢铁)
    (17.2746,交通运输)
    (17.7277,建筑材料)
    (18.4644,环保)
    (19.1048,公用事业)
    (19.3215,非银金融)
    (20.7985,石油石化)
    (22.7006,轻工制造)
    (24.2357,农林牧渔)
    (24.9682,医药生物)
    (25.2405,食品饮料)
    (28.9121,家用电器)
    (29.1692,基础化工)
    (29.4322,计算机)
    (30.7343,社会服务)
    (31.8911,电力设备)
    (34.6963,美容护理)
    (35.2568,汽车)
    (37.9843,商贸零售)
    (38.3215,机械设备)
    (38.7710,有色金属)
    (50.3083,电子)
    (52.7536,国防军工)
    (53.2198,综合)
    (57.5521,通信)
      };
      \addplot+[
        draw=none,
        fill={rgb,255:red,182;green,73;blue,73},
        bar shift=0pt
      ] coordinates {
        (-1.8940,房地产)
      };
      \draw[red!70!black,dashed,line width=0.8pt]
        (axis cs:28.1821,房地产) --
        (axis cs:28.1821,通信);
    \end{axis}
  \end{tikzpicture}
  \caption{2025 年 12 月 31 日申万一级行业原始 PE 中位数。样本包含
  5,438 条 PE 和行业分类均有效的股票记录；
  红色虚线为全市场原始 PE 中位数 28.18 倍，红色柱表示行业中位 PE 为负。}
  \label{fig:industry-raw-pe}
\end{figure}

行业中位 PE 最高的三个行业分别为通信（57.55 倍）、
综合（53.22 倍）和国防军工（52.75 倍）。
房地产的行业中位 PE 为 -1.89 倍，其负 PE 比例为
62.00\%；负 PE 表示行业内亏损观测较多，不应解释为负估值倍数。
本图仅为描述性统计，不包含回归检验。
